

\section{Data Terms}
\label{Sect_data_terms}

We assume that all data are trees over a fixed set of primitive
types. Most primitive types are standard.
We use 
$ {\bf char}, {\bf u64}, {\bf bool}, {\bf integer}, $ and
$ {\bf float}. $ Type $ {\bf u64} $ consists of 64-bit integers,
which can be used to index arrays. Type $ {\bf integer} $ 
consists of unbounded precision integers.  
In addition to the standard primitive types, there is a type
called $ {\bf selector}, $ which consists of named constants.
It is similar to an enumeration type, but with two differences:
Firstly, {\bf selectors} cannot be converted to (or obtained from) integers,
and secondly, there is only one $ {\bf selector} $
type. There is no need to allow
the user to define new enumeration types, because our type
system supports subtypes. Their intended use is to represent
tags that determine the types of later fields, 
like for examples the symbols in an attribute
grammar, or logical operators.
Compound trees are built using two constructors: The first
constructor constructs trees of fixed width. Once constructed,
the width cannot change.
The second constructor constructs trees of unbounded width,
with can grow and shrink like a vector. They are initially
constructed empty, after which elements can be inserted later.
\begin{definition}
   \label{Def_primtype} 
   We assume the following primitive types:
      {\bf bool}, {\bf char}, {\bf selector},
      {\bf u64}, {\bf integer}, {\bf double}.
   We call {\bf bool}, {\bf char}, and {\bf selector}
   \emph{tag types}, because they can be used as tags
   for selecting options.
   We assume that $ {\bf u64}, \ {\bf char} $ and $ {\bf bool} $
   are ordered types with the usual order.
\end{definition}

\begin{definition}
   \label{Def_data_term}
   We define the set of \emph{data terms} 
   by recursion as follows:
   \begin{itemize}
   \item
      If $ d $ is an element of one of the primitive types
      listed in Definition~\ref{Def_primtype},
      then $ d $ is a data term.
   \item
      If $ d_1, \ldots, d_n $ with $ n \geq 0 $ are
      data terms, then
      $ ( d_1, \ldots, d_n ) $
      and
      $ [ d_1, \ldots, d_n ] $ are also data terms. 
      We call $ ( d_1, \ldots, d_n ) $ a \emph{tuple}, and
      $ [ d_1, \ldots, d_n ] $ an \emph{array}.
   \end{itemize}
\end{definition}
The difference between $ (d_1, \ldots, d_n) $ and
$ [ d_1, \ldots, d_n ] $ is that the former represents
a fixed sized tuple, while the latter represents a vector.
In a program, the elements of a tuple $ (d_1, \ldots, d_n) $ can be
accessed by field access. For each field access instruction,
it is known which field it accesses.
At the same time, an array term $ [ d_1, \ldots, d_n ] $ can be indexed
by a variable of type $ {\bf u64}. $ 
\begin{definition}
   \label{Def_width}
   For a data term $ t, $ we define $ {\bf width}(t) $
   as the maximal size of a tuple occurring in $ t. $
   If $ t $ contains no tuples, then $ {\bf width}(t) = 0. $
\end{definition}
Because a program is always finite, 
there exists, for every program, an upperbound on the size
of the tuples that occur in it.
Hence, there always exists a natural number $ n, $ s.t.
all data terms that the program operates on, have
width at most $ n. $

\begin{example}
   Assume that we want to represent logical formulas  
   using operators $ \neg, \wedge, \vee, \rightarrow, \leftrightarrow, $ 
   where $ \neg $ is unary, $ \wedge $ and $ \vee $ are $ n $-ary, 
   and $ \leftrightarrow,\rightarrow $ are binary.
   For simplicity, we assume that atoms are represented
   by arrays of characters.

   The atom $ v0 $ can be represented by
   $ ( \sel{atom}, [ 'v', '0' ] ). $ 
   It is a tuple consisting of the selector constant
   \sel{atom} and an array, containing characters 'v' and '0'.
   
   The negation of $ v0 $ is $ ( \sel{neg}, v0 ). $ 
   Formula $ v1 \rightarrow v2 $ can be represented as
   $ ( \sel{implies}, ( \sel{atom}, [ 'v', '0' ] \, ), \,
                      ( \sel{atom}, [ 'v', '1' ] )). $ 
   Formula $ v_1 \wedge v_2 $ be be represented as
   $ ( \sel{and}, [ ( \sel{atom}, [ 'v', '0' ] \, ), \, 
                    ( \sel{atom}, [ 'v', '1' ] \, ) \, ] \, ). $ 
   The second element after an \sel{and}, or \sel{or} is always an
   array term. 
\end{example}
It can be seen from the example that, in constrast to most other
programming languages, logical operators are not represented by
constructors or function symbols, but by 
{\bf selector} constants. All formulas are constructed by
the same tuple constructor. For propositional formula,
the biggest arity that a tuple
can have, is three.
$ n $-ary formulas are represented by a binary tuple, whose
second element is an array.

